\documentclass[10pt,a3paper]{article}
\usepackage[margin=10mm]{geometry}
\usepackage[T1]{fontenc}
\usepackage[utf8]{inputenc}
\usepackage[french]{babel}
\usepackage{lmodern}
\usepackage{microtype}
\makeatletter\def\input@path{{../../../Style/}}\makeatother
\NeedsTeXFormat{LaTeX2e}
\ProvidesPackage{TLPosterCarteMentale}[2026/08/03 Posters et cartes mentales SI]

\RequirePackage{tikz}
\RequirePackage[most]{tcolorbox}
\RequirePackage{xcolor}
\RequirePackage{amsmath,amssymb}
\RequirePackage{graphicx}
\usetikzlibrary{positioning,arrows.meta,calc,shapes.geometric,mindmap,backgrounds,fit}

\definecolor{TLPosterBlue}{RGB}{24,86,141}
\definecolor{TLPosterLight}{RGB}{234,243,250}
\definecolor{TLPosterAccent}{RGB}{232,126,34}
\definecolor{TLPosterGreen}{RGB}{52,152,102}
\definecolor{TLPosterRed}{RGB}{192,57,43}
\definecolor{TLPosterGray}{RGB}{80,90,100}

\newcommand{\TLPosterTitre}[2]{%
  \begin{tcolorbox}[enhanced,boxrule=0pt,arc=0pt,colback=TLPosterBlue,left=6mm,right=6mm,top=4mm,bottom=4mm]
    {\color{white}\bfseries\LARGE #1}\hfill{\color{white}\large #2}
  \end{tcolorbox}%
}

\newtcolorbox{TLPosterBloc}[2][]{
  enhanced,breakable,
  colback=TLPosterLight,
  colframe=TLPosterBlue,
  colbacktitle=TLPosterBlue,
  coltitle=white,
  fonttitle=\bfseries,
  title={#2},
  boxrule=0.8pt,
  arc=2mm,
  left=3mm,right=3mm,top=2mm,bottom=2mm,
  #1
}

\newtcolorbox{TLPosterFormule}[1][]{
  enhanced,
  colback=white,
  colframe=TLPosterAccent,
  boxrule=1pt,
  arc=2mm,
  fontupper=\bfseries,
  halign=center,
  #1
}

\newcommand{\TLPosterMethode}[1]{%
  \begin{tcolorbox}[enhanced,colback=TLPosterGreen!8,colframe=TLPosterGreen,title=Méthode,fonttitle=\bfseries]
  #1
  \end{tcolorbox}%
}

\newcommand{\TLPosterAttention}[1]{%
  \begin{tcolorbox}[enhanced,colback=TLPosterRed!6,colframe=TLPosterRed,title=Point de vigilance,fonttitle=\bfseries]
  #1
  \end{tcolorbox}%
}

\tikzset{
  TLmindroot/.style={concept,concept color=TLPosterBlue,text=white,font=\bfseries\large,minimum size=34mm},
  TLmindlevel1/.style={level 1 concept/.append style={font=\bfseries,minimum size=23mm,sibling angle=45}},
  TLmindlevel2/.style={level 2 concept/.append style={font=\small,minimum size=17mm}},
  TLarrow/.style={-{Stealth[length=3mm]},thick,draw=TLPosterBlue},
  TLbox/.style={rounded corners=2mm,draw=TLPosterBlue,fill=TLPosterLight,align=center,inner sep=3mm}
}

\newenvironment{TLCarteMentale}[1]{%
  \begin{tikzpicture}[mindmap,grow cyclic,concept color=TLPosterBlue,every node/.style={concept},TLmindlevel1,TLmindlevel2]
  \node[TLmindroot]{#1}
}{;
  \end{tikzpicture}
}

\newcommand{\TLBranche}[3][]{child[concept color=#2]{node[#1]{#3}}}

\endinput

\pagestyle{empty}
\renewcommand{\familydefault}{\sfdefault}
\begin{document}
\TLPosterCourseHeader{3}{Lois entrée-sortie et transmetteurs}{Relier les paramètres d'entrée et de sortie d'un mécanisme}
\vspace{2mm}
\begin{minipage}[t]{0.61\linewidth}
\begin{TLPosterSavoirs}
\begin{itemize}
  \item distinguer coordonnées \textbf{articulaires} et coordonnées \textbf{opérationnelles}, modèle géométrique direct et inverse
  \item utiliser relation de Chasles, projections, norme et produit scalaire
  \item établir une fermeture géométrique pour une chaîne cinématique fermée
  \item obtenir une loi entrée-sortie en vitesse par dérivation des équations de position
  \item connaître les relations des transmetteurs usuels : engrenages, poulie--courroie, pignon--crémaillère, vis--écrou et trains épicycloïdaux
\end{itemize}
\end{TLPosterSavoirs}
\end{minipage}\hfill
\begin{minipage}[t]{0.36\linewidth}
\TLPosterKey{Quelle méthode ?}{Chaîne fermée : fermeture géométrique\\ Chaîne ouverte : modèle direct\\ Position imposée : modèle inverse\\ Vitesse : dériver la loi de position\\ Transmetteur : rapport cinématique}
\end{minipage}
\vfill
\begin{minipage}[t]{0.49\linewidth}
\begin{TLPosterBloc}[Chaîne fermée : loi en position]{TLPosterBlue}
Choisir les paramètres d'entrée/sortie puis écrire une boucle :
\[
\overrightarrow{AB}+\overrightarrow{BC}+\cdots+\overrightarrow{DA}=\vec 0.
\]
Projeter dans une base judicieuse et éliminer les paramètres intermédiaires. Une autre voie consiste à isoler un vecteur puis écrire sa \textbf{norme au carré} pour exploiter les produits scalaires.
\end{TLPosterBloc}
\begin{TLPosterBloc}[Chaîne ouverte : direct / inverse]{TLPosterPurple}
Le \textbf{modèle géométrique direct} donne la position de l'effecteur à partir des coordonnées articulaires. Pour imposer une position de l'effecteur, il faut résoudre le problème inverse : analytiquement lorsque c'est possible, sinon numériquement.
\end{TLPosterBloc}
\end{minipage}\hfill
\begin{minipage}[t]{0.49\linewidth}
\begin{TLPosterBloc}[Passer à la loi en vitesse]{TLPosterGreen}
Si la loi en position est connue, \textbf{dériver les équations scalaires} par rapport au temps : c'est souvent la méthode la plus rapide. Ne dériver qu'après avoir obtenu une relation géométrique correcte et conserver les variables intermédiaires tant qu'elles sont utiles.
\end{TLPosterBloc}
\begin{TLPosterBloc}[Transmetteurs : ne pas refaire une fermeture]{TLPosterTeal}
Pour un transmetteur usuel, utiliser directement sa relation cinématique et la convention de signe. Pour des engrenages extérieurs, le contact inverse le sens ; les rapports se multiplient dans une chaîne. Pour un train épicycloïdal, appliquer la relation de Willis. Vérifier ensuite vitesses, couple, puissance et rendement.
\end{TLPosterBloc}
\end{minipage}
\vfill
\begin{TLPosterMethode}
\begin{enumerate}
  \item Identifier clairement le \textbf{paramètre d'entrée} et le \textbf{paramètre de sortie} demandés.
  \item Reconnaître le cas : chaîne fermée, chaîne ouverte ou transmetteur usuel.
  \item Choisir la relation la plus courte : fermeture vectorielle, modèle direct ou rapport de transmission.
  \item Projeter ou utiliser une norme au carré pour éliminer les inconnues intermédiaires.
  \item Pour une loi en vitesse, dériver les équations de position obtenues.
  \item Tester la loi aux positions particulières et vérifier unités, signes, sens de mouvement et domaine géométrique.
\end{enumerate}
\end{TLPosterMethode}
\vfill
\TLPosterRetenir{ne pas calculer avant d'avoir identifié le type de mécanisme : fermeture, modèle direct/inverse et rapport de transmission répondent à des cas différents.}
\end{document}
